\documentclass[11pt,a4paper,openany]{book}

%==============================================================================
% COMBINED MONOGRAPH: Categorical Mechanics of Cytochrome P450
% Master document — includes the foundational partition-landscape paper
% and all fifteen papers of the Levinthal monograph series.
%==============================================================================

\usepackage[utf8]{inputenc}
\usepackage[T1]{fontenc}
\usepackage{amsmath,amssymb,amsthm,mathtools}
\usepackage{bm}
\usepackage{physics}
\usepackage{graphicx}
\usepackage{hyperref}
\usepackage[margin=2.5cm]{geometry}
\usepackage{enumitem}
\usepackage{float}
\usepackage{booktabs}
\usepackage{natbib}
\usepackage{xcolor}
\usepackage{caption}
\usepackage{microtype}
\usepackage{lmodern}
\usepackage[version=4]{mhchem}
\usepackage{multirow}
\usepackage{array}
\usepackage{pdfpages}   % for including compiled paper PDFs if available
\usepackage{tocloft}
\usepackage{titlesec}
\usepackage{fancyhdr}

% Theorem environments (shared)
\newtheorem{theorem}{Theorem}[chapter]
\newtheorem{lemma}[theorem]{Lemma}
\newtheorem{corollary}[theorem]{Corollary}
\newtheorem{proposition}[theorem]{Proposition}
\theoremstyle{definition}
\newtheorem{definition}[theorem]{Definition}
\newtheorem{axiom}{Axiom}[chapter]
\newtheorem{construction}[theorem]{Construction}
\theoremstyle{remark}
\newtheorem{remark}[theorem]{Remark}

% Custom commands
\newcommand{\kB}{k_{\mathrm{B}}}
\newcommand{\Depth}{\mathcal{M}}
\newcommand{\Part}{\mathcal{P}}
\newcommand{\dC}{d_{\mathrm{C}}}
\newcommand{\orderpar}{\langle r \rangle}
\newcommand{\kcat}{k_{\mathrm{cat}}}
\newcommand{\KM}{K_{\mathrm{M}}}
\newcommand{\nufloor}{\nu_{\mathrm{floor}}}
\newcommand{\Tpart}{T_{\Part}}
\newcommand{\Recv}{\mathcal{R}}

\hypersetup{
    colorlinks=true,
    linkcolor=blue,
    citecolor=blue,
    urlcolor=blue,
    pdftitle={Categorical Mechanics of Cytochrome P450},
    pdfauthor={Kundai Farai Sachikonye},
    pdfsubject={Partition Mechanics, Enzyme Catalysis, P450 Monograph}
}

% Headers
\pagestyle{fancy}
\fancyhf{}
\fancyhead[LE]{\small\textit{Categorical Mechanics of Cytochrome P450}}
\fancyhead[RO]{\small\textit{\rightmark}}
\fancyfoot[C]{\thepage}

%==============================================================================
\title{
    \vspace{2cm}
    {\LARGE\bfseries On the Geometric Consequences of Bounded Phase Space\\[0.4em]
    in Cytochrome P450}\\[1.2em]
    {\large\itshape On the Condition in Which Shakespeare Has Access to a Macbook,\\
    the Result Is Shakespeare}\\[1.5em]
    {\normalsize A Complete Monograph in Partition Dynamics}\\[2em]
    {\small Foundational Theory · Structural Construction · Observation ·
    Diversity · Consequences}
}
\author{
    \large Kundai Farai Sachikonye\\[0.5em]
    \normalsize AIMe Registry for Artificial Intelligence\\
    \normalsize Experimental Bioinformatics\\
    \normalsize Maximus-von-Imhof-Forum~3, 85354 Freising, Germany\\[0.5em]
    \normalsize\texttt{kundai.sachikonye@bitspark.com}
}
\date{\today}
%==============================================================================

\begin{document}

\frontmatter
\maketitle

%==============================================================================
% PREFACE
%==============================================================================
\chapter*{Preface}
\addcontentsline{toc}{chapter}{Preface}

This volume collects the complete Levinthal monograph on cytochrome P450,
together with its foundational theory paper.  The monograph was assembled in
the order in which the ideas were developed: first the general framework, then
its application to a specific enzyme family.

The foundational paper---\emph{On the Consequences of Global Partition
Structure Conservation: The Biological Partition Landscape}---established
that all biological dynamics follows from a single geometric principle:
bounded phase space admits partition and nesting.  That paper proved the claim
across five biological domains with 39 validation tests and $89\,\%$ overall
accuracy.

The fifteen papers of the P450 monograph then applied this principle to a
single enzyme family at a depth that the foundational paper could not reach.
They constitute a bottom-up derivation of everything that is known about
cytochrome P450: from the abstract observation that any biological molecule
must have both a receiver and a sender, through sequence-space geometry,
through the complete catalytic cycle, through 57 human isoforms, through
spectroscopy and database recovery, to a predictive model of pharmacogenomics
under polypharmacy.

The prefatory chapter (Chapter~\ref{ch:intro}) introduces the framework and
provides a guided tour of the monograph.  Readers already familiar with the
partition-landscape formalism may proceed directly to Part~II.

\bigskip\noindent
\textit{Kundai Farai Sachikonye}\\
Freising, \today

%==============================================================================
\tableofcontents
%==============================================================================

\mainmatter

%==============================================================================
% CHAPTER 0: INTRODUCTION / STAGE-SETTER
%==============================================================================

\chapter{Introduction and Framework}
\label{ch:intro}
\markboth{Introduction and Framework}{Introduction and Framework}

%------------------------------------------------------------
\section{On the Disrespect of Shakespeare}
\label{sec:shakespeare}
%------------------------------------------------------------

Levinthal's paradox, as it is taught, comes with a popular illustration: a monkey at a typewriter, given infinite time, eventually produces \textit{Hamlet}.  The implication is that the protein's success is the kind of vanishingly rare event that random sampling could in principle achieve.  The framing has shaped how three generations of biochemists have thought about folding.  The framing is wrong on every count, and the wrongness is computable.

The monkey-Shakespeare comparison treats Shakespeare's output as a sequence of characters that some procedure could reproduce.  No procedure can.  Not the monkey.  Not the human typist.  Not Shakespeare himself.  \textit{Hamlet} is what one specific evaluation, in one specific receiver, at one specific point in the receiver's trajectory, returned.  The output exists; the evaluation is closed; nothing produces \textit{Hamlet} twice.  The protein is in the same position.  Levinthal's $3^N$ conformational space is structurally identical to the monkey's $47^{130{,}000}$ key-space: a calculation of how large a search would have to be \emph{if search were what was happening}.  Search is not what is happening.  The protein evaluates its receiver.  The native fold is what the evaluation returns.  For CYP3A4 at 503 residues, $3^{503} \approx 10^{240}$; the framework completes the fold in six categorical steps.  The reduction is not a clever algorithm; it is the dissolution of the search frame.

The full argument---quantifying Shakespeare's stylometric signature, running the monkey calculation ($10^{140}$ total keystrokes for all atoms in the observable universe at one key per Planck time, still short by 216{,}860 orders of magnitude), establishing the category error (the monkey is not random; it is doing what monkeys do), the Newton corollary (a Macbook does not help Newton write \textit{Principia} because the bottleneck is the receiver, not the instrument), and the forward-is-rewind point (even Shakespeare cannot produce Shakespeare twice, because the partition count has incremented)---is given in full in the corresponding prefatory paper.

\textit{What would Shakespeare do with a Macbook?  Type like Shakespeare.  The Macbook just sits there.  What does cytochrome~P450 do with bounded phase space?  It evaluates its receiver.  The seventeen papers that follow exhibit what that evaluation returns.}

%------------------------------------------------------------
\section{The Problem of Fragmentation}
\label{sec:problem}
%------------------------------------------------------------

Biology is described by a museum of theories.  Transition-state theory
quantifies activation barriers~\citep{eyring1935}.  Flux-balance analysis
redistributes metabolic flow~\citep{orth2010}.  Marcus theory addresses
electron tunnelling.  Pharmacokinetic models handle drug clearance.  Sequence
homology explains isoform diversity.  Population genetics predicts
polymorphism phenotypes.  Each framework introduces its own vocabulary,
parameters, and limits.  None derives from the others.

This monograph argues that the fragmentation reflects a missing foundation,
not a fundamental heterogeneity.  The missing foundation is
\emph{partition structure}.

All biological systems share one property: they are \emph{bounded}.  Cells
occupy finite volumes.  Proteins assume finite conformations.  Reactions occur
in finite time.  Electron transfer crosses finite distances.  This boundedness
is not a constraint imposed from outside; it is the defining character of a
localized physical entity.  Bounded systems admit partition and nesting.
Partition depth $\Depth$ measures how much hierarchical structure is required
to distinguish a state from the totality of alternatives.  Every dynamical law
follows as a corollary of this single quantity.

The biological-partition-landscape paper~\citep{sachikonye2025_bpl} proved
these claims across five domains---atomic structure, enzyme kinetics,
metabolism, disease, and cellular organisation---achieving $89\,\%$ predictive
accuracy with zero free parameters.  This monograph executes the same
programme for a single family---cytochrome P450---at a depth that the broader
paper could not reach.  The fifteen papers constitute a bottom-up synthesis:
from the initial observation that biological molecules must have both a
receiver and a sender, through sequence-space geometry, through single-chain
folding, through cofactor positioning and equilibrium states, through the
complete catalytic cycle, to a universal classification of all~57 human
isoforms, all reaction types, all spectroscopic observables, and all
pharmacogenomic outcomes.  Everything emerges from~$\Depth$.

%------------------------------------------------------------
\section{Foundations: Bounded Phase Space and Partition Depth}
\label{sec:intro-foundations}
%------------------------------------------------------------

\subsection{Boundedness and Finite Partition}

\begin{axiom}[Boundedness]
\label{ax:bounded}
Every localized physical system occupies a bounded region of phase space
$\Omega \subset \mathbb{R}^{2n}$ with $\mathrm{Vol}(\Omega) < \infty$.
\end{axiom}

\begin{axiom}[Finite Partition]
\label{ax:partition}
A bounded phase space admits partition into a finite number of distinguishable
cells,
\begin{equation}
    \Omega = \bigsqcup_{i=1}^{N} \Omega_i, \qquad
    N = \frac{\mathrm{Vol}(\Omega)}{h^n} < \infty,
\end{equation}
where $h^n$ is the minimum cell volume set by the uncertainty principle.
\end{axiom}

These two axioms are observations, not postulates.  The uncertainty relation
$\Delta x\,\Delta p \geq \hbar/2$ establishes a hard floor on cell volume;
finite total volume then guarantees finite cell count.  Partitions can nest:
cells contain subcells, which contain sub-subcells, forming a hierarchy.

\subsection{Partition Coordinates and Atomic Structure}

The hierarchical structure of nested partitions defines a coordinate system.
For atomic systems, the coordinates are
\[
    n \geq 1\;(\text{nesting depth}),\quad
    \ell \in \{0,\ldots,n{-}1\}\;(\text{complexity}),\quad
    m \in \{-\ell,\ldots,+\ell\}\;(\text{orientation}),\quad
    s \in \{\pm\tfrac{1}{2}\}\;(\text{parity}).
\]
The number of distinct states at depth $n$ is $C(n) = 2n^2$.  This formula
exactly reproduces atomic shell structure---$C(1)=2$, $C(2)=8$, $C(3)=18$,
$C(4)=32$---without postulation.  The partition coordinate system \emph{is}
the quantum number system; both arise from the same geometric necessity.

\subsection{Partition Depth}

\begin{definition}[Partition Depth]
The partition depth $\Depth$ of a categorical state is the cumulative
logarithmic branching required for its unique specification:
\begin{equation}
    \Depth = \sum_{i} \log_b(k_i),
    \label{eq:depth}
\end{equation}
where $k_i$ is the branching factor at hierarchy level $i$ and $b$ is the
encoding base.
\end{definition}

$\Depth$ measures how much a state must be distinguished from the totality of
alternatives.  The connection to thermodynamic entropy is exact:

\begin{theorem}[Depth--Entropy Equivalence]
$S_{\Part} = \kB\,\Depth\,\ln b$.
\end{theorem}

Partition depth is not invented for biology; it is the quantity that already
accounts for the periodic table, the structure of information channels, and the
cost of biological specificity.

%------------------------------------------------------------
\section{Dynamical Laws from Partition Geometry}
\label{sec:intro-dynamics}
%------------------------------------------------------------

\subsection{The Existence Cost}

Before any process can occur, its participants must \emph{exist} as distinct
entities---they must be partitioned from the undifferentiated totality of
reality.  This has an unavoidable cost.

\begin{theorem}[Existence Cost]
For any process involving entities $\{A_i\}$,
\begin{equation}
    E_{\mathrm{exist}} = \kB\,\Tpart\,\ln b \cdot \sum_i \Depth_{A_i},
    \label{eq:exist}
\end{equation}
where $\Tpart = 65$~kJ~mol$^{-1}$ is the characteristic partition temperature.
\end{theorem}

\subsection{Partition Time and Gradient Flow}

By the energy-time uncertainty relation, a partition operation with energy
$\sim\kB T$ requires time
\begin{equation}
    \tau_p = \frac{\hbar}{\kB T} \approx 25~\text{fs at }298~\text{K}.
    \label{eq:tauP}
\end{equation}

Since partitioning is not instantaneous, and reality does not pause, every
transition leaves an irreducible residue---entropy production is built in at
the foundation level.

\begin{theorem}[Partition Gradient Flow]
\begin{equation}
    \frac{d\mathbf{x}}{dt} = -\gamma\,\nabla\Depth(\mathbf{x}).
\end{equation}
Systems evolve toward states of lower partition depth.  What we call forces are
partition depth gradients; there are no forces in a deeper sense.
\end{theorem}

\subsection{Transition States and the Universal Rate Law}

The transition state between A and B must encode where the system came from,
where it is going, and which path connects them:
\[
    \Depth_T = \Depth_A + \Depth_B + \Depth_{\mathrm{path}}.
\]

The activation depth is $\Delta\Depth = \Depth_T - \Depth_A$.  In condensed
phase, the effective floor frequency is
\begin{equation}
    \nufloor = \kappa_{\mathrm{sol}} \cdot \frac{1}{\tau_p}
             = 10^{10}~\text{s}^{-1},
    \label{eq:nufloor}
\end{equation}
where $\kappa_{\mathrm{sol}} \approx 2.5\times10^{-4}$ captures solvent
friction, recrossing, and tunnelling corrections in water at physiological
temperature.  The rate law is:

\begin{equation}
    \boxed{k = \nufloor\,\exp(-\Delta\Depth)}
    \label{eq:rate}
\end{equation}

This is the Arrhenius equation derived from geometry.  $\nufloor$ and
$\Tpart$ are not fitted to P450 data; they are set once from the partition-
landscape calibration and held fixed throughout.

%------------------------------------------------------------
\section{Catalysis as Partition Topology Provision}
\label{sec:intro-catalysis}
%------------------------------------------------------------

\subsection{What an Enzyme Does}

An enzyme provides an alternative partition topology with reduced transition
depth:
\begin{equation}
    \Depth_T^{\mathrm{cat}} = \Depth_T^{\mathrm{uncat}} - \Depth_{\mathrm{shared}}.
\end{equation}
It does not change $\Depth_A$ or $\Depth_B$ (equilibria are unaffected); it
reduces the barrier by lending its own partition structure to the transition
state.

\subsection{The Ball Game}

The physics is most transparent through the ball-game analogy of
Ref.~\citep{sachikonye2025_bpl}.  Two teams on opposite sides of a partitioned
wall score by throwing balls through holes.  What matters is proximity to a
hole, not ball speed.  Each score forces rearrangement (entropy generation).
The game cannot reverse.

In this picture:
\begin{itemize}[nosep]
    \item An enzyme adds or enlarges holes (reduces $\Delta\Depth$).
    \item A competitive inhibitor occupies holes without scoring.
    \item A mechanism-based inactivator welds holes shut.
    \item An inducer builds new walls with more holes.
    \item A polymorphism narrows or blocks existing holes.
\end{itemize}
Every pharmacological intervention in Part~III of this monograph is a specific
modification of the hole geometry.

\subsection{Categorical Distance and Efficiency}

\begin{definition}[Categorical Distance]
$\dC(A, B)$ is the number of partition boundaries that must be crossed in going
from state A to state B.
\end{definition}

\begin{theorem}[Efficiency from Categorical Distance]
\begin{equation}
    \frac{\kcat}{\KM} \propto \frac{1}{\dC \cdot \tau_p}
    \quad\Longrightarrow\quad
    \log_{10}\!\left(\frac{\kcat}{\KM}\right) \approx 10 - \dC.
    \label{eq:efficiency}
\end{equation}
Each additional partition boundary costs approximately one order of magnitude
in catalytic efficiency.
\end{theorem}

The four-hop NADPH$\to$FAD$\to$FMN$\to$heme chain of cytochrome P450
reductase has $\dC = 4$, predicting $\kcat/\KM \sim 10^6$~M$^{-1}$s$^{-1}$;
the experimentally measured CPR--P450 coupling value matches this within
error.

%------------------------------------------------------------
\section{Why Cytochrome P450?}
\label{sec:why-p450}
%------------------------------------------------------------

A framework that claims universality should face the hardest test available.
Cytochrome P450 provides that test for three reasons.

\textbf{Breadth.}  The CYP superfamily comprises 57 human isoforms covering a
substrate space of tens of thousands of small molecules.  Any framework that
predicts only one isoform is limited; this monograph covers all 57 as points
on a single sequence manifold parameterised by $\Depth$.

\textbf{Mechanistic depth.}  The P450 catalytic cycle spans seven distinct
states, an obligate four-electron oxygen activation, heterolytic O--O cleavage
to Compound~I (a formal Fe$^{\mathrm{IV}}$=O porphyrin $\pi$-cation radical),
and a C--H activation step with KIE up to 10.  Each step is a
partition-topology transition; each is modelled at the spectroscopic level by
an independent instrument stack.

\textbf{Pharmacological consequence.}  CYP3A4 alone clears $\sim50\,\%$ of
all clinical drugs.  CYP2D6 phenotypes span a seven-fold rate range that
determines codeine toxicity.  CYP2C9~$\ast$3 reduces S-warfarin hydroxylation
to $<5\,\%$ of wild-type.  The framework must reproduce compound phenotypes
under polypharmacy---a much stronger test than predicting a single
$k_{\mathrm{cat}}$.

No other enzyme family offers this combination of mechanistic richness,
structural diversity, and clinical relevance.

%------------------------------------------------------------
\section{Notation and Constants}
\label{sec:notation}
%------------------------------------------------------------

Two empirical constants are fixed from the partition-landscape calibration:

\medskip
\begin{center}
\begin{tabular}{lll}
\toprule
Symbol & Value & Meaning \\
\midrule
$\nufloor$ & $10^{10}$~s$^{-1}$ & Condensed-phase floor frequency \\
$\Tpart$   & $65$~kJ~mol$^{-1}$ & Partition temperature per depth unit \\
\bottomrule
\end{tabular}
\end{center}

\medskip
All rates, $K_i$ values, $\Delta M$ shifts, and compound phenotype
predictions in the fifteen papers derive from these two numbers.  They are
not re-fitted to any paper.

The rate law $k = \nufloor\,\exp(-\Delta\Depth)$ and the efficiency relation
$\log_{10}(\kcat/\KM) \approx 10 - \dC$ are the two equations that govern
the entire monograph.  All other formulae are applications or corollaries.

The additive compound-phenotype formula (Paper~15) is:
\begin{equation}
    \Delta\Depth_{\mathrm{eff}}
      = \Delta\Depth_{\mathrm{allele}}
      + \sum_{j} \Delta\Delta\Depth_j,
    \label{eq:phenotype}
\end{equation}
where each $\Delta\Delta\Depth_j = \ln\alpha_j$ with
$\alpha_j = 1 + [I_j]/K_{i,j}$ is the inhibitor depth shift for the $j$-th
co-administered drug.

%------------------------------------------------------------
\section{Roadmap of the Monograph}
\label{sec:roadmap}
%------------------------------------------------------------

\begin{description}[style=nextline]

\item[Part~I: Foundational Theory (Chapter~\ref{ch:bpl})]
    \emph{On the Consequences of Global Partition Structure Conservation:
    The Biological Partition Landscape.}  The generative paper---proves the
    two axioms and derives all corollaries across five biological domains.

\item[Part~II: Construction (Papers 1--3 and 2.5)]
    Expression algebra and the receiver (P1); the sequence manifold (P2);
    GLB structural input (P2.5); CYP3A4 equilibrium states (P3).

\item[Part~III: Observation (Papers 4--6)]
    Five-layer electron transfer observation (P4, the headline);
    Compound~I formation (P5); C--H activation and rebound (P6).

\item[Part~IV: Diversity and Reactions (Papers 7--10)]
    Heteroatom oxidation and dealkylation (P7);
    atypical reactions atlas (P8);
    57-isoform taxonomy (P9);
    pharmacogenomics atlas (P10).

\item[Part~V: Physical Machinery and Synthesis (Papers 11--12)]
    Membrane anchoring and CPR coupling (P11);
    seven-state closed orbit (P12).

\item[Part~VI: Spectroscopy and Informatics (Papers 13--14)]
    Spectroscopic atlas (P13);
    ternary database recovery (P14).

\item[Part~VII: Pharmacology (Paper 15)]
    Polymorphisms, DDI, and compound phenotype prediction (P15).
    132 validation scripts: all PASS.

\end{description}

%==============================================================================
% PART I: FOUNDATIONAL THEORY
%==============================================================================

\part{Foundational Theory}
\label{part:foundation}

%==============================================================================
\chapter{The Biological Partition Landscape}
\label{ch:bpl}
\markboth{The Biological Partition Landscape}{The Biological Partition Landscape}
%==============================================================================

\begin{center}
\textit{On the Consequences of Global Partition Structure Conservation:\\
The Biological Partition Landscape}\\[0.5em]
\small Kundai Farai Sachikonye\\
Technical University of Munich, School of Life Sciences
\end{center}

\bigskip

\noindent\textbf{Abstract.} We derive a theory of biological dynamics from a
single geometric principle: physical systems occupy bounded regions of phase
space that admit partition and nesting.  From this principle, we establish
partition depth $\Depth$ as the fundamental quantity governing all processes.
We prove: (i)~any process requires partitioning participants from the totality
of reality, incurring an unavoidable existence cost; (ii)~dynamics flow along
partition depth gradients; (iii)~catalysts function by providing alternative
partition topologies with reduced depth barriers; (iv)~living systems are
self-maintaining partition structures.  Applied to enzyme catalysis, we derive
reaction rates, substrate specificity, and catalytic efficiency from partition
geometry alone.  Applied to metabolism, we derive pathway architecture and
flux distributions.  Applied to disease, we derive pathology as partition
malformation.  All results emerge as corollaries of the central theorem with
zero free parameters.  We validate predictions against enzyme kinetic data
spanning six orders of magnitude in catalytic efficiency, metabolic flux
measurements, and disease progression correlations.

\bigskip

\textit{The full paper is in}
\texttt{publication/biological-partition-landscape/biological-partition-landscape.tex}

%% If the BPL paper has been compiled, uncomment the following line:
% \includepdf[pages=-]{../../publication/biological-partition-landscape/biological-partition-landscape.pdf}

%==============================================================================
% PART II: CONSTRUCTION
%==============================================================================

\part{Construction}
\label{part:construction}

%==============================================================================
\chapter{Paper 1: Expression Algebra for Biomolecules}
\label{ch:p1}
%==============================================================================
\textit{Full paper:}
\texttt{cytochrome/publications/foundations/expression-algebra-for-biomolecules/}
\texttt{expression-algebra-proteins.tex}
% \includepdf[pages=-]{../cytochrome/publications/foundations/expression-algebra-for-biomolecules/expression-algebra-proteins.pdf}

%==============================================================================
\chapter{Paper 2: The P450 Address Manifold}
\label{ch:p2}
%==============================================================================
\textit{Full paper:}
\texttt{cytochrome/publications/foundations/p450-sequence-space/}
\texttt{cytochrome-p450-sequence-space.tex}
% \includepdf[pages=-]{../cytochrome/publications/foundations/p450-sequence-space/cytochrome-p450-sequence-space.pdf}

%==============================================================================
\chapter{Paper 2.5: GLB Structural Input}
\label{ch:p25}
%==============================================================================
\textit{Full paper:}
\texttt{cytochrome/publications/foundations/glb-structural-input/}
\texttt{glb-structural-input.tex}
% \includepdf[pages=-]{../cytochrome/publications/foundations/glb-structural-input/glb-structural-input.pdf}

%==============================================================================
\chapter{Paper 3: CYP3A4 Equilibrium States}
\label{ch:p3}
%==============================================================================
\textit{Full paper:}
\texttt{cytochrome/publications/equilibrium-states/cyp3a4-resting-substrate-bound/}
\texttt{cyp3a4-resting-substrate-bound.tex}
% \includepdf[pages=-]{../cytochrome/publications/equilibrium-states/cyp3a4-resting-substrate-bound/cyp3a4-resting-substrate-bound.pdf}

%==============================================================================
% PART III: OBSERVATION
%==============================================================================

\part{Observation}
\label{part:observation}

%==============================================================================
\chapter{Paper 4: Electron Transfer Chain (Headline)}
\label{ch:p4}
%==============================================================================
\textit{Full paper:}
\texttt{cytochrome/publications/catalytic-cycle/multi-hop-et-chain/}
\texttt{multi-hop-et-chain.tex}
% \includepdf[pages=-]{../cytochrome/publications/catalytic-cycle/multi-hop-et-chain/multi-hop-et-chain.pdf}

%==============================================================================
\chapter{Paper 5: Compound I Formation}
\label{ch:p5}
%==============================================================================
\textit{Full paper:}
\texttt{cytochrome/publications/catalytic-cycle/compound-i-formation/}
\texttt{compound-i-formation.tex}
% \includepdf[pages=-]{../cytochrome/publications/catalytic-cycle/compound-i-formation/compound-i-formation.pdf}

%==============================================================================
\chapter{Paper 6: C--H Activation and Rebound}
\label{ch:p6}
%==============================================================================
\textit{Full paper:}
\texttt{cytochrome/publications/catalytic-cycle/ch-activation-rebound/}
\texttt{ch-activation-rebound.tex}
% \includepdf[pages=-]{../cytochrome/publications/catalytic-cycle/ch-activation-rebound/ch-activation-rebound.pdf}

%==============================================================================
% PART IV: DIVERSITY AND REACTIONS
%==============================================================================

\part{Diversity and Reactions}
\label{part:diversity}

%==============================================================================
\chapter{Paper 7: Heteroatom Oxidation and Dealkylation}
\label{ch:p7}
%==============================================================================
\textit{Full paper:}
\texttt{cytochrome/publications/reactions/heteroatom-dealkylation/}
\texttt{heteroatom-dealkylation.tex}
% \includepdf[pages=-]{../cytochrome/publications/reactions/heteroatom-dealkylation/heteroatom-dealkylation.pdf}

%==============================================================================
\chapter{Paper 8: Atypical Reactions Atlas}
\label{ch:p8}
%==============================================================================
\textit{Full paper:}
\texttt{cytochrome/publications/reactions/atypical-reactions-atlas/}
\texttt{atypical-reactions-atlas.tex}
% \includepdf[pages=-]{../cytochrome/publications/reactions/atypical-reactions-atlas/atypical-reactions-atlas.pdf}

%==============================================================================
\chapter{Paper 9: 57-Isoform Taxonomy}
\label{ch:p9}
%==============================================================================
\textit{Full paper:}
\texttt{cytochrome/publications/diversity/57-human-isoforms/}
\texttt{57-human-isoforms.tex}
% \includepdf[pages=-]{../cytochrome/publications/diversity/57-human-isoforms/57-human-isoforms.pdf}

%==============================================================================
\chapter{Paper 10: Pharmacogenomics Atlas}
\label{ch:p10}
%==============================================================================
\textit{Full paper:}
\texttt{cytochrome/publications/pharmacology/pharmacogenomics-atlas/}
\texttt{pharmacogenomics-atlas.tex}
% \includepdf[pages=-]{../cytochrome/publications/pharmacology/pharmacogenomics-atlas/pharmacogenomics-atlas.pdf}

%==============================================================================
% PART V: PHYSICAL MACHINERY AND SYNTHESIS
%==============================================================================

\part{Physical Machinery and Synthesis}
\label{part:machinery}

%==============================================================================
\chapter{Paper 11: Membrane Anchoring and CPR Coupling}
\label{ch:p11}
%==============================================================================
\textit{Full paper:}
\texttt{cytochrome/publications/construction/membrane-cofactor-cpr/}
\texttt{membrane-cpr.tex}
% \includepdf[pages=-]{../cytochrome/publications/construction/membrane-cofactor-cpr/membrane-cpr.pdf}

%==============================================================================
\chapter{Paper 12: Seven-State Closed Orbit}
\label{ch:p12}
%==============================================================================
\textit{Full paper:}
\texttt{cytochrome/publications/synthesis/seven-state-closed-orbit/}
\texttt{closed-orbit.tex}
% \includepdf[pages=-]{../cytochrome/publications/synthesis/seven-state-closed-orbit/closed-orbit.pdf}

%==============================================================================
% PART VI: SPECTROSCOPY AND INFORMATICS
%==============================================================================

\part{Spectroscopy and Informatics}
\label{part:spectroscopy}

%==============================================================================
\chapter{Paper 13: Spectroscopic Atlas}
\label{ch:p13}
%==============================================================================
\textit{Full paper:}
\texttt{cytochrome/publications/spectroscopy/spectroscopic-atlas/}
\texttt{spectroscopic-atlas.tex}
% \includepdf[pages=-]{../cytochrome/publications/spectroscopy/spectroscopic-atlas/spectroscopic-atlas.pdf}

%==============================================================================
\chapter{Paper 14: Ternary Database Recovery}
\label{ch:p14}
%==============================================================================
\textit{Full paper:}
\texttt{cytochrome/publications/informatics/database-recovery/}
\texttt{database-recovery.tex}
% \includepdf[pages=-]{../cytochrome/publications/informatics/database-recovery/database-recovery.pdf}

%==============================================================================
% PART VII: PHARMACOLOGY
%==============================================================================

\part{Pharmacology}
\label{part:pharmacology}

%==============================================================================
\chapter{Paper 15: Polymorphisms, DDI, and Compound Phenotype}
\label{ch:p15}
%==============================================================================
\textit{Full paper:}
\texttt{cytochrome/publications/diversity/polymorphisms-ddi-inhibitors/}
\texttt{polymorphisms-ddi-inhibitors.tex}
% \includepdf[pages=-]{../cytochrome/publications/diversity/polymorphisms-ddi-inhibitors/polymorphisms-ddi-inhibitors.pdf}

%==============================================================================
% BACK MATTER
%==============================================================================

\backmatter

%==============================================================================
\chapter*{Validation Summary}
\addcontentsline{toc}{chapter}{Validation Summary}
%==============================================================================

\begin{center}
\begin{tabular}{lrrl}
\toprule
Paper & Scripts & Status & Domain \\
\midrule
Partition Landscape (BPL) & 39 tests & 89\,\% pass & Five domains \\
\midrule
Paper 1  & 8  & All PASS & Expression algebra \\
Paper 2  & 8  & All PASS & Sequence manifold \\
Paper 2.5 & 8 & All PASS & Structural input \\
Paper 3  & 8  & All PASS & Equilibrium states \\
Paper 4  & 8  & All PASS & ET chain / apparatus \\
Paper 5  & 8  & All PASS & Compound I \\
Paper 6  & 8  & All PASS & C--H activation \\
Paper 7  & 8  & All PASS & Heteroatom reactions \\
Paper 8  & 8  & All PASS & Atypical reactions \\
Paper 9  & 8  & All PASS & 57-isoform taxonomy \\
Paper 10 & 8  & All PASS & Pharmacogenomics \\
Paper 11 & 8  & All PASS & Membrane / CPR \\
Paper 12 & 8  & All PASS & Closed orbit \\
Paper 13 & 8  & All PASS & Spectroscopic atlas \\
Paper 14 & 8  & All PASS & Database recovery \\
Paper 15 & 8  & All PASS & Polymorphisms / DDI \\
\midrule
\textbf{Total} & \textbf{132} & \textbf{All PASS} & \\
\bottomrule
\end{tabular}
\end{center}

%==============================================================================
\chapter*{Consolidated References}
\addcontentsline{toc}{chapter}{Consolidated References}
%==============================================================================

\bibliographystyle{apalike}
\bibliography{%
    ../cytochrome/publications/introduction/references,%
    ../cytochrome/publications/foundations/expression-algebra-for-biomolecules/references,%
    ../cytochrome/publications/catalytic-cycle/multi-hop-et-chain/references,%
    ../cytochrome/publications/diversity/polymorphisms-ddi-inhibitors/references%
}

\end{document}
